% !TeX program = xelatex
\documentclass[aspectratio=169,14pt]{beamer}
\usetheme{Serula}

\title[Serula]{Serula}
\subtitle{The slides are further along}
\author{Your name}
\institute{Department of unfinished business}
\date{}

\begin{document}

\begin{frame}[t,noframenumbering]
  \titlepage
\end{frame}

\section{Structure}
\begin{frame}[t]{Lists and small digressions}
  \framesubtitle{Bullets, numbers, and a block}
  \begin{columns}[T,onlytextwidth]
    \begin{column}{0.48\textwidth}
      \subheading{Unordered list}
      \begin{itemize}
        \item A point.
          \begin{itemize}
            \item A supporting detail.
              \begin{itemize}
                \item A small digression.
              \end{itemize}
          \end{itemize}
        \item Back to the point.
      \end{itemize}
    \end{column}
    \begin{column}{0.48\textwidth}
      \subheading{Ordered list}
      \begin{enumerate}
        \item Duplicate a slide.
          \begin{enumerate}
            \item Change the text.
              \begin{enumerate}
                \item This bit too.
              \end{enumerate}
          \end{enumerate}
        \item Read it once out loud.
      \end{enumerate}
    \end{column}
  \end{columns}
  \vspace{1mm}
  \begin{block}{A note on space}
    You can leave some of it empty.
  \end{block}
\end{frame}

\section{Numbers}
\begin{frame}[t]{Equations and tables}
  \framesubtitle{One rule, three examples}
  \begin{columns}[T,onlytextwidth]
    \begin{column}{0.53\textwidth}
      {\small A number multiplied by itself:\par}
      \begin{align*}
        f(x) &= x^2,\\
        f(3) &= 9.
      \end{align*}
    \end{column}
    \begin{column}{0.40\textwidth}
      \begin{table}
        \centering
        \small
        \renewcommand{\arraystretch}{1.2}
        \begin{tabular}{@{}rr@{}}
          \toprule
          Number & Square \\
          \midrule
          1 & 1 \\
          2 & 4 \\
          3 & 9 \\
          \bottomrule
        \end{tabular}
        \caption{Three examples.}
      \end{table}
    \end{column}
  \end{columns}
  \vspace{1mm}
  {\small\itshape The next row is left to the reader.\par}
\end{frame}

\newcommand{\orbitdiagram}[1]{%
  \begin{tikzpicture}[x=1cm,y=1cm]
    \usebeamercolor{structure}
    \path[use as bounding box] (-2.5,-1.3) rectangle (2.5,1.3);
    \fill[structure.fg!6] (0,0) circle[radius=1.05];
    \draw[structure.fg!45,line width=0.8pt] (0,0) circle[radius=1.05];
    \draw[Secondary!50,densely dashed] (0,0) -- (#1:1.05);
    \fill[Secondary] (0,0) circle[radius=0.035];
    \fill[structure.fg] (#1:1.05) circle[radius=0.12];
  \end{tikzpicture}%
}

\section{Visuals}
\begin{frame}[t]{Figures}
  \framesubtitle{A drawing and an animation}
  \begin{columns}[T,onlytextwidth]
    \begin{column}{0.47\textwidth}
      \begin{figure}
        \centering
        \orbitdiagram{0}
        \caption{A dot with nowhere urgent to be.}
      \end{figure}
    \end{column}
    \begin{column}{0.47\textwidth}
      \begin{figure}
        \centering
        \begin{animateinline}[
          controls,loop,poster=first,buttonsize=3mm,
          buttonfg=0.2:0.2:0.7,
          alttext={A dot moving around a circular path}
        ]{12}
          \multiframe{36}{iAngle=0+10}{\orbitdiagram{\iAngle}}
        \end{animateinline}
        \caption{The same dot, taking the long way.}
      \end{figure}
    \end{column}
  \end{columns}
  \vspace{2mm}
  {\small\color{Secondary}Use the playback controls in a compatible PDF viewer.\par}
\end{frame}

\section{References}
\begin{frame}[t]{References}
  \framesubtitle{Here go your sources}
  \begin{thebibliography}{9}
    \bibitem{sample-book}
      Author name.
      \newblock \emph{Book title}.
      \newblock Publisher, year.

    \bibitem{sample-article}
      Author name and Coauthor name.
      \newblock Article title.
      \newblock \emph{Journal name}, volume(issue):pages, year.
  \end{thebibliography}
\end{frame}

\end{document}